% 事件运算的集合示意：并、交、补、互斥与完备事件组。纯示意图，无数据。
\begin{tikzpicture}[line width=.6pt,font=\fontsize{8}{10}\selectfont,>=Stealth]
\def\W{2.8}\def\H{2.0}\def\R{0.62}
\tikzset{omega/.style={draw=FigureStroke,line width=.6pt},
  ring/.style={draw=FigureBlue,line width=.75pt}}
% (a) 并
\begin{scope}[xshift=0cm]
 \fill[FigureBlueFill] (1.05,1.0) circle[radius=\R] (1.75,1.0) circle[radius=\R];
 \draw[ring] (1.05,1.0) circle[radius=\R] (1.75,1.0) circle[radius=\R];
 \draw[omega] (0,0) rectangle (\W,\H);
 \node[FigureInk] at (0.75,1.0) {$A$}; \node[FigureInk] at (2.05,1.0) {$B$};
 \node[FigureInk,anchor=north west] at (0.02,\H) {$\Omega$};
 \node[FigureInk] at (\W/2,-0.35) {(a) $A\cup B$：至少一个发生};
\end{scope}
% (b) 交
\begin{scope}[xshift=3.12cm]
 \begin{scope}\clip (1.05,1.0) circle[radius=\R]; \fill[FigureBlueFill] (1.75,1.0) circle[radius=\R]; \end{scope}
 \draw[ring] (1.05,1.0) circle[radius=\R] (1.75,1.0) circle[radius=\R];
 \draw[omega] (0,0) rectangle (\W,\H);
 \node[FigureInk] at (0.75,1.0) {$A$}; \node[FigureInk] at (2.05,1.0) {$B$};
 \node[FigureInk,anchor=north west] at (0.02,\H) {$\Omega$};
 \node[FigureInk] at (\W/2,-0.35) {(b) $A\cap B$：同时发生};
\end{scope}
% (c) 补
\begin{scope}[xshift=6.24cm]
 \fill[FigureBlueFill,even odd rule] (0,0) rectangle (\W,\H) (1.4,1.0) circle[radius=\R];
 \draw[ring] (1.4,1.0) circle[radius=\R];
 \draw[omega] (0,0) rectangle (\W,\H);
 \node[FigureInk] at (1.4,1.0) {$A$};
 \node[FigureInk,anchor=north west] at (0.02,\H) {$\Omega$};
 \node[FigureInk] at (\W/2,-0.35) {(c) $A^{\mathsf c}$：$A$不发生};
\end{scope}
% (d) 互斥
\begin{scope}[xshift=9.36cm]
 \fill[FigureBlueFill] (0.8,1.0) circle[radius=0.5];
 \fill[FigureGreenFill] (2.0,1.0) circle[radius=0.5];
 \draw[ring] (0.8,1.0) circle[radius=0.5];
 \draw[draw=FigureGreen,line width=.75pt] (2.0,1.0) circle[radius=0.5];
 \draw[omega] (0,0) rectangle (\W,\H);
 \node[FigureInk] at (0.8,1.0) {$A$}; \node[FigureInk] at (2.0,1.0) {$B$};
 \node[FigureInk,anchor=north west] at (0.02,\H) {$\Omega$};
 \node[FigureInk] at (\W/2,-0.35) {(d) 互斥：$A\cap B=\varnothing$};
\end{scope}
% (e) 完备事件组
\begin{scope}[xshift=12.48cm]
 \fill[FigureBlueFill] (0,0) rectangle (0.9,\H);
 \fill[FigureGreenFill] (0.9,0) -- (2.0,0) -- (1.7,\H) -- (0.9,\H) -- cycle;
 \fill[FigureYellowFill] (2.0,0) -- (\W,0) -- (\W,\H) -- (1.7,\H) -- cycle;
 \draw[FigureStroke,line width=.6pt] (0.9,0)--(0.9,\H) (2.0,0)--(1.7,\H);
 \draw[omega] (0,0) rectangle (\W,\H);
 \node[FigureInk] at (0.45,0.9) {$B_1$}; \node[FigureInk] at (1.4,0.9) {$B_2$}; \node[FigureInk] at (2.3,0.9) {$B_3$};
 \node[FigureInk,anchor=north west] at (0.02,\H) {$\Omega$};
 \node[FigureInk] at (\W/2,-0.35) {(e) 完备事件组};
\end{scope}
\end{tikzpicture}
